\section{Experimental design}
\label{sec:design}

\subsection{Two dataset builds and two runs}
Two training runs used the identical model, preprocessing and optimisation settings, one per build (Section~\ref{sec:builds}): run~1 on build~1 (at most eight phase-2 epochs) and run~2 on build~2 (at most fourteen phase-2 epochs and a wall-clock budget). Both added sources of build~2 were part of the original data design with labelling rules fixed before training, but had not finished downloading when build~1 was made. Adding them was a data intervention that used neither the test split nor the external cohort for any parameter. However, run~2 was started after run~1's test and external results had been seen, and those results bear on the question that build~2 addresses. The external cohort therefore cannot be regarded as untouched by the research process. Run~1's external results were produced at about 09:28 on 2026-11-03 and run~2 was started at about 10:07 the same day. A strictly prospective evaluation of the selected model needs a new external cohort.

\subsection{Run selection rule}
\label{sec:selection}
Because the two validation splits differ in composition, neither test nor external results may be used to choose between runs. A rule was therefore fixed in advance and recorded in the repository before run~2's results existed. The primary criterion is macro-F1 on a \emph{common selection set} of validation images that neither build used for training or testing; the higher value is selected, and if the difference is below 0.01, run~1 is kept unless run~2's macro ROC-AUC is higher by at least 0.01. Accuracy, macro ROC-AUC and the same metrics restricted to the four sources present in both builds are reported alongside, because run~2 has an in-distribution advantage on the TBX11K and COVID-19 Radiography images. An exact McNemar comparison of paired predictions is reported only as discordant-pair counts. The run not selected is a secondary experiment.

\subsection{Follow-up runs}
\label{sec:followup}
Three further runs on the selected build were fixed before run~2 had been evaluated (Table~S5 in the supplement): runs 3 and 4 change only the seed (43 and 44) and measure seed variability, and run~5 keeps the seed and removes the three filters (same manifest and split), as a preprocessing ablation. Each uses the phase-2 epoch limit that run~2 actually completed, with early stopping (patience 3). They are trained sequentially because parallel training does not fit in 8~GB of memory. Seed variability is the mean, sample standard deviation, minimum and maximum over runs 2--4; the ablation is one run read against that spread. The served model remains run~2, chosen by the rule above; choosing the best of several seeds on test or external results would be selection on the evaluation sets.

\subsection{Once-only evaluation, analyses and reproducibility}
For each run, the test split and the external cohort were evaluated once after the checkpoint had been fixed on validation, and no setting was changed afterwards. The analyses are overall and per-class metrics, per-source accuracy, calibration, the external binary task, the run comparison, the source probe, seed variability, the ablation, Grad-CAM, and the C1--C4 analyses with the C2 decision.

Training and evaluation ran on a laptop CPU with 8~GB of memory under Windows (Python~3.12, TensorFlow~2.16, Keras~3, no GPU). Runs are therefore single seeds (\RSeed) or three-seed groups, and each is bounded by an epoch limit or a wall-clock budget; the reported run trained for \REpochsTrained{} epochs and its selected checkpoint is from epoch \RSelectedEpoch. Every run writes a metadata record (code version, configuration snapshot, manifest SHA-256, split fingerprint, seed, library versions, per-epoch history, class weights, checkpoint SHA-256), and every evaluation report repeats the checkpoint SHA-256 (\texttt{\RCheckpointShort\ldots}, version \RModelVersion), manifest SHA-256 and split fingerprint. All numbers in this paper are generated from those reports into a macro file. The images cannot be redistributed under the source terms; the repository contains the rules, builder, provenance logs and checksums needed to rebuild the identical dataset from the official sources.
